% !TEX root = ../b-rg.tex
\chapter{Pronouns}\label{ch:pronouns}

Pronouns are the one part of the Borlish nominal system that still distinguishes case. The personal pronouns have separate subject, accusative, oblique and disjunctive forms, though no person keeps all four apart, and the third person singular distinguishes gender, which nouns do not. This chapter describes the personal and reflexive pronouns, the possessives, the agentive pronouns that are characteristic of Borlish relative clauses, and the indefinite and interrogative pronouns. The demonstrative pronouns are treated with the demonstrative determiners in Chapter~\ref{ch:determiners}.

\section{Personal pronouns}\label{sec:personal-pronouns}
The personal pronouns are given in the table below.
\begin{table}[ht]
    \centering
    \begin{tabular}{l l l l l l}
        \multicolumn{6}{c}{Personal pronouns} \\
        \hline
        & subject & accusative & oblique & disjunctive & reflexive \\
        \hline
        1\textsc{sg} & \bl{jo} & \bl{me} & \bl{me} & \bl{mey} & \bl{me} \\
        2\textsc{sg} & \bl{tu} & \bl{te} & \bl{te} & \bl{tey} & \bl{te} \\
        3\textsc{sg.m} & \bl{le} & \bl{le} & \bl{luy} & \bl{luy} & \bl{se} \\
        3\textsc{sg.f} & \bl{la} & \bl{la} & \bl{luy} & \bl{luy} & \bl{se} \\
        3\textsc{sg.n} & \bl{lo} & \bl{lo} & \bl{luy} & \bl{luy} & \bl{se} \\
        1\textsc{pl} & \bl{nos} & \bl{nos} & \bl{nos} & \bl{nos} & \bl{se} \\
        2\textsc{pl} & \bl{vos} & \bl{vos} & \bl{vos} & \bl{vos} & \bl{se} \\
        3\textsc{pl} & \bl{ðe} & \bl{lou} & \bl{lour} & \bl{lou} & \bl{se} \\
        \hline
    \end{tabular}
\end{table}
No person keeps all four cases apart. In the first and second person singular, accusative and oblique share a form (\bl{me}, \bl{te}). In the third person singular, the accusative is identical to the subject form, and a single oblique and disjunctive form \bl{luy} serves all three genders. The first and second person plural have a single invariable form each. Only the third person plural distinguishes three forms: \bl{ðe}, \bl{lou} and \bl{lour}.

Gender is distinguished only in the third person singular: \bl{le} `he', \bl{la} `she' and \bl{lo} `it'. Since nouns have no grammatical gender (Chapter~\ref{ch:nouns}), the choice is made by meaning: \bl{le} and \bl{la} for male and female persons, and \bl{lo} for everything else, including animals of unspecified sex, objects, abstractions and whole propositions. The second person plural \bl{vos} is also used as a polite singular (Chapter~\ref{ch:register-names}).

Several forms from different parts of the paradigm are homophones: \bl{tu}, \bl{tey} and the possessive \bl{ty} are all \phm{ti}; \bl{mey} and \bl{my} are \phm{mi}; \bl{sey} and \bl{sy} are \phm{si}.
    \glphm{jo tu le la lo nos vos ðe | me te luy lou lour | mey tey sey}{ʒo ti le la lo nɔz vɔz ðe | me te laj lu lur | mi ti si}{{\First\Sg} {\Second\Sg} {\Third\Sg.\M} {\Third\Sg.\F} {\Third\Sg.\N} {\First\Pl} {\Second\Pl} {\Third\Pl} | {\First\Sg.\Obl} {\Second\Sg.\Obl} {\Third\Sg.\Obl} {\Third\Pl.\Acc} {\Third\Pl.\Obl} | {\First\Sg.\Dsj} {\Second\Sg.\Dsj} {\Refl.\Dsj}}

\section{Subject pronouns}\label{sec:subject-pronouns}

\subsection{Expression of the subject}
Borlish does not drop its subjects: a finite verb has an expressed subject even where its ending would identify the person. The only exceptions are imperatives and the expletive `it' (see below).
    \glsen{La me dau ðorn tranq ag massour.}{la me do ðɔrn trank ɛj maˈsur}{{\Third\Sg.\F} {\First\Sg.\Obl} {give-\Pst} punch precise {at.\Def} jaw}{She punched me right in the jaw.}
A few verbs regularly place their subject after the verb. The most important are verbs of need and fitness such as \bl{stovir} `be needed, be wanted' and \bl{valir} `be worth, ought'. Their experiencer is an oblique pronoun before the verb, and their subject, the thing needed or wanted, follows the verb and controls its agreement (Chapter~\ref{ch:voice}). The subject may be a noun phrase or an infinitive clause.
    \glsen{Nos stoun balis sobottal luicessem.}{nɔz stun baˈlɪz ˌso.bɔˈtal lajˈtsɛ.sɛm}{{\First\Pl} {be.needed-\Third\Pl} beacon radio {bright-\Cmpr}}{We need stronger radio transmitters.}
    \glsen{Me stou for savir ig nos seyau scaðr den.}{me stu fɔr saˈvɪr aj nɔz siˈjo ˈxa.ðr̩ dɛn}{{\First\Sg.\Obl} {be.wanted} only {know-\Inf} {\Comp} {\First\Pl} {be-\First\Pl.\Sbjv} aligned thereof}{I want to know we're on the same page.}
Several copular verbs likewise take a following subject, among them \bl{cair} `fall', which in this construction means `happen, come about'.
    \glsen{Capitan, cay uncos inovrant cas y motour prophthan.}{ˌka.piˈtan ke ɪnˈkɔz ˌi.nɔˈvrant kaz i moˈtur prɔfˈθan}{captain fall something {malfunction-\Prs.\Ptcp} regarding {\Def} engine {faster.than.light}}{Captain, something's gone wrong with the warp drive.}
The expletive subject `it' is usually \bl{lo}, as in \bl{Parcanç es lo temerer\ldots} `It is perhaps foolhardy\ldots', but it is often simply left out: \bl{Parcanç sera megl vos ignoreð\ldots} `Perhaps it is better you not know\ldots'.

\subsection{Position}
A subject pronoun normally stands before the finite verb, separated from it only by object clitics and preverbal negation. In questions, and when another constituent is placed first in the clause, it follows the finite verb instead (Chapter~\ref{ch:simple-clause}).
    \glsen{Nuvr so jo my seucq historic molinant.}{ˈnivr̩ so ʒo mi sawk ˌhɪ.stoˈrɪk ˌmo.liˈnant}{recently {be.\First\Sg} {\First\Sg} {\First\Sg.\Poss} path historic {mill-\Prs.\Ptcp}}{Lately I have been contemplating my legacy.}
    \glsen{Posc cay prinvir porreu nos alcheðr.}{pɔx ke prɪnˈvir pɔˈraw nɔz alˈke.ðr̩}{after fall spring {can-\Fut.\First\Pl} {\First\Pl} anew}{Next spring we can start over.}
    \glsen{Son ðe nervous d'issaur ny cors?}{sɔn ðe nɛrˈvuz diˈsor ni kɔrz}{{be.\Third\Pl} {\Third\Pl} confident {of=great.success} {in.\Def} contest}{Are they sure they'll win the contest?}

\subsection{Elision and contraction}
Before a word beginning with a vowel, \bl{jo} is elided to \bl{j'}, and \bl{le}, \bl{la} and \bl{lo} all to \bl{l'}, so that the elided form no longer shows gender. The third person plural \bl{ðe} is never elided.
    \glsen{J'am çamoc sciftar ag tostuð.}{ʒam tsaˈmɔk xɪfˈtar ɛj tɔˈstɪθ}{{\First\Sg=like} jam {spread-\Inf} {at.\Def} toast}{I like to have jam on toast.}
    \glsen{Ðe ayen heu lastumnesc, com ig ðe no'rn for pog dormið.}{ðe ɛˈjɛn haw ˌla.stɪmˈnɛx kɔm aj ðe nɔrn fɔr pɔj dɔrˈmɪθ}{{\Third\Pl} {have-\Ipfv-\Third\Pl} look weary as {\Comp} {\Third\Pl} {\Neg=\Cop.\Ipfv-\Third\Pl} only little {sleep-\Pst.\Ptcp}}{They looked haggard, as though they hadn't slept much.}
With the copula, the subject pronouns instead form contractions in which the gender of the third person is kept: \bl{jo's} \phm{ʒɔz} `I am', \bl{tu's} \phm{tɪz} `you are', and \bl{le's, la's, lo's} \phm{lɛz laz lɔz} `he, she, it is'. In the past, \bl{le'r, la'r, lo'r} \phm{lɛr lar lɔr} `he, she, it was' contract the subject with \bl{er} (Chapter~\ref{ch:copulas}).
    \glsen{J'osc sy cigl novel; la's mallovicher sad demay.}{ʒox si tsajl noˈvɛl laz ˌma.lo.vɪˈkɛr sad deˈme}{{\First\Sg=hate} {\Third\Sg.\Poss} sweetheart new {\Third\Sg.\F=be} pert gravely {too.much}}{I hate her new girlfriend; she's far too chipper.}
    \glsen{Le'r mell'yam sleghant avant my rau!}{lɛr mɛˈljam sliˈjant aˈvant mi ro}{{\Third\Sg.\M=\Cop.\Ipfv} {\First\Sg.\Poss-\Def=uncle} {slander-\Prs.\Ptcp} {in.front.of} {\First\Sg.\Poss} disdain}{He was shamelessly slandering my uncle!}

\section{Object pronouns}\label{sec:object-pronouns}

\subsection{Position}
The accusative and oblique pronouns are clitics. They stand immediately before the verb form whose object they are: a finite verb, an infinitive or a participle. Since objects in Borlish precede non-finite verbs (Chapter~\ref{ch:simple-clause}), a pronoun object of an infinitive or participle follows the auxiliary or modal and precedes the non-finite form. It does not climb to the modal.
    \glsen{Ðe la macquaurn vars surcair un tarç dagr.}{ðe la maˈkworn varz sɪrˈker ɪn tarts dɛjr}{{\Third\Pl} {\Third\Sg.\F.\Acc} {raise-\Pst.\Third\Pl} towards {end.up-\Inf} {\Indf} third secret}{They raised her to be neither fish nor fowl.}
    \glsen{Jo no poð lou affaçar pre vienn y flottisc.}{ʒo no pɔθ lu ˌa.faˈtsar pre ˈvjɛnn̩ i flɔˈtɪx}{{\First\Sg} {\Neg} can {\Third\Pl.\Acc} {face-\Inf} before {come-\Third\Pl} {\Def} fleet}{I can't confront them until reinforcements arrive.}
    \glsen{J'ay un desc farrago nos cogt.}{ʒe ɪn dɛx ˌfa.raˈgo nɔz kɔjt}{{\First\Sg=have.\First\Sg} {\Indf} meal hash {\First\Pl} {cook.\Pst.\Ptcp}}{I've cooked us some hash.}

\subsection{Combinations and elision}
When two clitics occur together, a first- or second-person pronoun comes before a third-person one, as in \bl{Jo te lo dis} `I told you so'. Among third-person clitics, the reflexive \bl{se} comes first, then the accusative, then the oblique: \bl{se lo}, \bl{lo luy}. Before a vowel, \bl{me} and \bl{te} are elided to \bl{m'} and \bl{t'}, and \bl{le, la, lo} to \bl{l'}.
    \glsen{Jo m'asseyeu ny caðer magrastr cos coðelous.}{ʒo ˌma.siˈjau ni kaˈðɛr mɛjˈrastr̩ kɔz ˌko.ðeˈluz}{{\First\Sg} {\First\Sg.\Obl=seat-\Pst} {in.\Def} chair rickety {\Adv} wary}{I sat gingerly on the rickety chair.}
    \glsen{T'incombisce te leyar tostessem.}{tɪnˌkɔm.biˈxe te liˈjar tɔˈstɛ.sɛm}{{\Second\Sg.\Obl=behove-\Ipfv} {\Second\Sg.\Obl} {raise-\Inf} {early-\Cmpr}}{You should have got up earlier.}

\subsection{Imperatives}
In commands, an object pronoun follows the verb and takes its disjunctive form. This holds for every command, whether expressed with the imperative or the subjunctive, positive or negative.
    \glsen{Dig mey; tu va noc a y los me çaveyar.}{daj mi ti va nɔk a i lɔz me ˌtsa.viˈjar}{say {\First\Sg.\Dsj} {\Second\Sg} go {\Neg} to {\Def} slab {\First\Sg.\Obl} {cushion-\Inf}}{Tell me; you don't have to sugarcoat it.}
    \glsen{Indigneð mey noc dell'oc baugl van.}{ˌɪn.dajˈnɛθ mi nɔk dɛˈlɔk bojl van}{{begrudge-\Second\Pl.\Sbjv} {\First\Sg.\Dsj} {\Neg} {of.\Def=\Prox.\Sg} camp vain}{Don't hold this bad set-up against me.}

\section{Disjunctive pronouns}\label{sec:disjunctive}
Besides following imperatives, the disjunctive pronouns are used after prepositions, and in other positions where the pronoun is not attached to a verb: in coordination, in isolation, and when the pronoun is emphasised. A coordinated subject containing a disjunctive takes plural agreement on the verb.
    \glsen{Mey eð y maðr jognau l'oç kivougl tras y ralley.}{mi ɛθ i ˈmaðr̩ ʒɔjˈno lɔts kiˈvujl traz i raˈli}{{\First\Sg.\Dsj} and {\Def} mother {join-\Ipfv.\First\Pl} {\Def=\Dist.\Sg} jigsaw across {\Def} afternoon}{Me and my mum did the jigsaw all afternoon.}
    \glsen{J'er pol frigsant je ados cas mey.}{ʒɛr pɔl frajˈzant ʒe aˈdɔz kaz mi}{{\First\Sg=\Cop.\Ipfv} chicken {fry-\Prs.\Ptcp} {day.of} behind {at.home.of} {\First\Sg.\Dsj}}{I was frying chicken at home yesterday.}
    \glsen{Nos erau runam alleyant con lou commet a fin.}{nɔz eˈro riˈnam ˌa.liˈjant kɔn lu kɔˈmɛt a fɪn}{{\First\Pl} {\Cop.\Ipfv-\First\Pl} secretly {ally-\Prs.\Ptcp} with {\Third\Pl.\Dsj} start to end}{We were secretly on their side the whole time.}
A disjunctive pronoun may be followed by \bl{meðes} `self' for emphasis, as in \bl{mey meðes} `myself' or \bl{avant luy meðes} `in front of him(self)'. After any other constituent, \bl{meðes} means `even'.
    \glsen{La pegntau toð y lettranç ag reðrpuy meðes.}{la pijnˈto tɔθ i lɛˈtrants ɛj ˌrɛ.ðr̩ˈpaj meˈðɛz}{{\Third\Sg.\F} {paint-\Pst} all {\Def} text {at.\Def} background even}{She even painted all the background text.}

\section{Reflexive and reciprocal pronouns}\label{sec:reflexive}
The reflexive pronoun is \bl{se}, elided to \bl{s'} before a vowel, with the disjunctive \bl{sey}. In the first and second person singular, the ordinary object forms \bl{me} and \bl{te} are used reflexively.
    \glsen{Jo me landau y visaç raut nell'evier.}{ʒo me lanˈdo i viˈzats rot ˌnɛ.leˈvjɛr}{{\First\Sg} {\First\Sg.\Obl} {wash-\Pst} {\Def} face quick {in.\Def=sink}}{I quickly washed my face in the sink.}
    \glsen{Le s'yauvnau y cavel sborreyað cos van.}{le sjovˈno i kaˈvɛl ˌsbɔ.riˈjaθ kɔz van}{{\Third\Sg.\M} {\Refl=pat-\Pst} {\Def} hair unkempt {\Adv} vain}{He failed to pat down his messy hair.}
In the plural, \bl{se} is used with all three persons, including first and second person subjects.
    \glsen{Nos se contentaum a un tatover soul porçonnar.}{nɔz se ˌkɔn.tɛnˈtom a ɪn ˌta.toˈvɛr sul ˌpɔr.tsɔˈnar}{{\First\Pl} {\Refl} {content-\Pst-\First\Pl} to {\Indf} potato alone {share-\Inf}}{We made do with sharing a single potato.}

Beyond true reflexives, \bl{se} has several other uses (see also Chapter~\ref{ch:voice}). It marks reciprocal action, sometimes reinforced by a disjunctive with \bl{meðes}, as in \bl{se cavir lou meðes a sgart} `to see each other'.
    \glsen{L'ec fraðr narc se feirn acconnosc eld veint catr.}{lɛk fraðr̩ nark se firn̩ ˌa.kɔˈnɔx ɛld vint katr̩}{{\Def=\Prox.\Pl} brother step {\Refl} {do.\Pst-\Third\Pl} acquaintance {at.age} twenty four}{The stepbrothers met when they were twenty-four.}
It forms intransitive counterparts of transitive verbs, as in \bl{Y lamp se quïsceu} `The light went out', and, with \bl{cavir} `get' and a past participle, a passive.
    \glsen{L'accommission se cof interront par y sounç d'enmigtað.}{ˌla.kɔ.mɪˈsjɔn se kɔf ˌɪn.tɛˈrɔnt par i sunts ˌdɛn.majˈtaθ}{{\Def=undertaking} {\Refl} {get.\Pst} {interrupt.\Pst.\Ptcp} by {\Def} release {of=enmity}}{The project was interrupted by the outbreak of hostilities.}
Finally, with \bl{poðir} `can' and other verbs, \bl{se} serves as an impersonal subject `one, you'. Borlish has no other generic pronoun.
    \glsen{Se poðe pom accatar teint ogtoç solt a cascun.}{se poˈðe pɔm ˌa.kaˈtar tint ɔjˈtɔts sɔlt a kaˈxɪn}{{\Refl} {can-\Ipfv} apple {buy-\Inf} {amount-\Prs.\Ptcp} eighteen penny at {each.one}}{You could buy apples for eighteen pence each.}

\section{Possessives}\label{sec:possessives}
Each person has a possessive determiner, used before a noun, and a possessive pronoun, used independently.
\begin{table}[ht]
    \centering
    \begin{tabular}{l l l l}
        \multicolumn{4}{c}{Possessives} \\
        \hline
        & determiner & before a vowel & pronoun \\
        \hline
        1\textsc{sg} & \bl{my} & \bl{mell'} & \bl{mien} \\
        2\textsc{sg} & \bl{ty} & \bl{tell'} & \bl{tien} \\
        3\textsc{sg} & \bl{sy} & \bl{sell'} & \bl{sien} \\
        1\textsc{pl} & \bl{nossy} & \bl{nossell'} & \bl{nostr} \\
        2\textsc{pl} & \bl{vossy} & \bl{vossell'} & \bl{vostr} \\
        3\textsc{pl} & \bl{lorry} & \bl{lorrell'} & \bl{lour} \\
        \hline
    \end{tabular}
\end{table}
The possessive determiners take no article before a consonant, and fuse with it before a vowel (Section~\ref{sec:def-article}). The third person singular \bl{sy} does not distinguish gender: it means `his', `her' or `its' according to the possessor.
    \glsen{Mell'amig lou narreyau vars agir temerer.}{ˌmɛ.laˈmaj lu ˌna.riˈjo varz aˈʒɪr ˌte.meˈrɛr}{{\First\Sg.\Poss-\Def=friend} {\Third\Pl.\Acc} {goad-\Pst} towards {act-\Inf} rash}{My friend goaded them into acting rashly.}
The possessive pronouns take the definite article.
    \glsen{Ty jahesment es angostessem sur y mien.}{ti ˌʒa.hɛˈzmɛnt ɛz ˌan.gɔˈstɛ.sɛm sɪr i mjɛn}{{\Second\Sg.\Poss} set be {thin-\Cmpr} on {\Def} {\First\Sg.\Poss}}{Your TV set is thinner than mine.}
In literary style, a possessive pronoun may instead follow the noun, as in \bl{toð y hayennaç sien} `all her workload'. Outside literature this survives in fixed expressions such as \bl{Deu mien!} `My God!' and \bl{thein vostr} `at your service'. It is less marked after an indefinite, as in \bl{un amig mien} `a friend of mine', and some speakers use it freely with other determiners, preferring \bl{l'oc amig mien} to \bl{mell'oc amig} `this friend of mine'.

\section{Agentive pronouns}\label{sec:agentive}
Borlish has a set of agentive pronouns, meaning `by me, by you' and so on: \bl{meyon} \phm{miˈjɔn}, \bl{teyon} \phm{tiˈjɔn}, \bl{noscon} \phm{noˈxɔn}, \bl{voscon} \phm{voˈxɔn}, and \bl{seyon} \phm{siˈjɔn}, which serves the third person in both numbers. They express the agent of a past participle or of a potential form in \bl{-abr}, and stand immediately before it. Their main use is in the participial relative clauses that are the commonest way of relativising on an object in Borlish (Chapter~\ref{ch:relative-clauses}).
    \glsen{My fraðr dessoclau y castel blader meyon costroit.}{mi ˈfra.ðr̩ ˌde.soˈklo i kaˈstɛl blaˈdɛr miˈjɔn koˈstrɔjt}{{\First\Sg.\Poss} brother {topple-\Pst} {\Def} castle {card-\Adj} {\First\Sg.\Agt} {build.\Pst.\Ptcp}}{My brother knocked over the house of cards I built.}
    \glsen{Nos nolau hur dar for an svattour noscon confiabr.}{nɔz noˈlo hɪr dar fɔr an svaˈtur noˈxɔn kɔnˈfja.br̩}{{\First\Pl} {\Neg.\Fut-\Ipfv.\First\Pl} hire {give.\Inf} only {to.\Indf} detective {\First\Pl.\Agt} {trust-\Pot}}{We were only going to hire a PI we could trust.}
An object clitic may stand between the agentive pronoun and the participle.
    \glsen{Y sovragl seyon me dað er adornað par fylau.}{i soˈvrɛjl siˈjɔn me daθ ɛr ˌa.dɔrˈnaθ par fiˈlo}{{\Def} belt {\Third.\Agt} {\First\Sg.\Obl} {give.\Pst.\Ptcp} {\Cop.\Ipfv} {decorate-\Pst.\Ptcp} by crow}{The belt she gave me was decorated with crows.}
The agentive pronouns have no noun counterpart. Where the agent is a full noun phrase, it follows the participle in a prepositional phrase with \bl{de} or \bl{par}.

\section{Indefinite and negative pronouns}\label{sec:indefinite-pronouns}
The principal indefinite pronouns are \bl{un hom} `someone, anyone' (literally `a person') and \bl{uncos} `something, anything'.
    \glsen{M'orruð un hom ay my bos oundrað.}{mɔˈrɪθ ɪn ɔm e mi bɔz unˈdraθ}{{\First\Sg.\Obl=seem} {\Indf} person {have.\Sbjv} {\First\Sg.\Poss} gun {wound-\Pst.\Ptcp}}{I think someone's broken my gun.}
    \glsen{Vos vouð uncos special accommettr por y gaviscenç ou?}{vɔz vuθ ɪnˈkɔz speˈtsjal ˌa.kɔˈmɛtr̩ pɔr i ˌga.viˈxɛnts u}{{\Second\Pl} {will.\Second\Pl} something special {undertake-\Inf} for {\Def} celebration {\Q}}{Are you going to do anything special for the festivities?}
The negative counterparts are \bl{alcun hom} `nobody' and \bl{nien} `nothing'; \bl{necun} `neither, none' is used of a choice among several, as in \bl{necun dy conjuraçonnour} `none of the conspirators'. On their combination with negation, see Chapter~\ref{ch:negation}.
    \glsen{Alcun hom au y lamp fornið con zygom.}{alˈkɪn ɔm o i lamp fɔrˈnɪθ kɔn ziˈgɔm}{no person {have-\Pst} {\Def} light {provide.\Pst.\Ptcp} with switch}{No-one had installed a switch for the light.}
    \glsen{Ðe n'aun nien ostendað tras dou meis.}{ðe non njɛn ˌɔ.stɛnˈdaθ traz du miz}{{\Third\Pl} {\Neg=have.\Third\Pl} nothing {post-\Pst.\Ptcp} across two month}{They've posted nothing in two months.}
The universal pronouns are \bl{toð gent} and \bl{casc hom} `everyone', \bl{cascun} `each one' and \bl{tout} `everything, all'. \bl{Tout} can also float after the verb, referring to a plural subject.
    \glsen{Nos coum tout eurað par un grand jolleistr.}{nɔz kum tut awˈraθ par ɪn grant ʒoˈlistr̩}{{\First\Pl} {get.\Pst-\First\Pl} all drunk by {\Indf} big {night.out}}{We all got drunk during a big night out.}
Free-choice `any(thing), whatever' is \bl{calscon}, as in \bl{calscon fallent} `whatever is needed', and `whoever' is \bl{cuscon}, as in \bl{cuscon sey le} `whoever he may be'.

\section{Interrogative and relative pronouns}\label{sec:interrogative-pronouns}
The interrogative pronouns are \bl{cu} \phm{ki} `who' and \bl{cal} `what, which', which is also the interrogative determiner (Section~\ref{sec:quantifiers}). Both are used in direct and indirect questions (Chapter~\ref{ch:sentence-types}).
    \glsen{Cu a l'oc ugnlac vorpal laiscað jos ty panner?}{ki a lɔk ajnˈlak vɔrˈpal leˈxaθ ʒɔjz ti paˈnɛr}{who have {\Def=\Prox.\Sg} onion diced {leave-\Pst.\Ptcp} {next.to} {\Second\Sg.\Poss} breadbin}{Who's left this diced onion next to the breadbin?}
    \glsen{Jo lo memour noc cu protagoney l'oç lasc.}{ʒo lo meˈmur nɔk ki proˌta.goˈni lɔts lax}{{\First\Sg} {\Third\Sg.\N.\Acc} remember {\Neg} who {star.in} {\Def=\Dist.\Sg} film}{I can't remember who stars in that film.}
\bl{Cuscon} `whoever' is also used as an interrogative `who', as in \bl{Cuscon dirigeu l'administraçon\ldots?} `Who led the government\ldots?'.
% CHECK: how do interrogative cu and cuscon differ? Register, or a nuance such
% as 'who on earth'?

The relative pronoun \bl{dont} `of which, of whom, whose' introduces finite relative clauses, as in \bl{ogt ivan, dont Adagle sta tressem} `eight children, of whom Adaillé was the third'. Free relatives `what, that which' are formed with the demonstrative \bl{l'oç} (Section~\ref{sec:dem-without-art}). Relative clauses are treated fully in Chapter~\ref{ch:relative-clauses}.
